\begin{beginnerstatement}[Kurz erklärt: Qualität und Nachweise]

Qualitätssicherung umfasst die Maßnahmen, mit denen Fehler gesucht und
Anforderungen überprüft werden. Ein Test ist eine konkret ausgeführte Prüfung.
Verifikation bedeutet hier, nachvollziehbar zu kontrollieren, ob ein erwartetes
technisches Verhalten tatsächlich belegt ist. Evidenz oder ein Nachweis ist
das prüfbare Ergebnis; ein Artefakt kann dabei beispielsweise ein erzeugtes
Paket oder ein Protokoll sein. Ein Commit ist ein eindeutig gespeicherter
Quellcodestand, weshalb ein commitgebundener Nachweis nur für genau diesen
Stand gilt. \emph{PASS} steht für bestanden, \emph{FAIL} für ausgeführt und
fehlgeschlagen und \emph{SKIP} für nicht ausgeführt oder nicht anwendbar.
Fehlende Evidenz heißt deshalb \emph{nicht verifiziert} und ist weder ein
Erfolg noch automatisch ein Fehler.

\end{beginnerstatement}
